\documentclass[aspectratio=169,11pt]{beamer}
\usepackage[utf8]{inputenc}
\usepackage[T1]{fontenc}
\usepackage{lmodern}
\usepackage[english]{babel}
\usepackage{graphicx}
\usepackage{booktabs}
\usepackage{amsmath,amssymb}
\usepackage{xcolor}
\usepackage{listings}
\usetheme{Madrid}
\usecolortheme{default}
\setbeamertemplate{navigation symbols}{}
\setbeamertemplate{footline}[frame number]
\setbeamertemplate{blocks}[rounded][shadow=false]
\AtBeginSection[]{\begin{frame}{Outline}\tableofcontents[currentsection,subsectionstyle=hide]\end{frame}}
\definecolor{hsbiblue}{RGB}{45,57,120}
\definecolor{hsbimid}{RGB}{76,114,176}
\definecolor{hsbilight}{RGB}{225,231,247}
\definecolor{cgreen}{RGB}{85,164,103}
\definecolor{cred}{RGB}{196,78,82}
\lstset{basicstyle=\ttfamily\scriptsize,breaklines=true,keywordstyle=\color{hsbiblue},commentstyle=\color{cgreen},columns=fullflexible,showstringspaces=false}
\title{Fast Track --- Mock Exam 2 (Solutions)}
\subtitle{ShopSphere scenario $\cdot$ full worked solutions}
\author{Prof.\ Dr.\ Christoph Weisser}
\institute{HSBI}
\date{Summer Semester 2026}
\begin{document}
\begin{frame}\titlepage\end{frame}

\begin{frame}{How to use these solutions}
\begin{itemize}
  \item These slides give a \textbf{fully worked answer} to every question of Mock Exam 2 (the ShopSphere paper), one question per slide.
  \item For \textbf{multiple-choice} questions the correct option is stated, followed by \emph{why each distractor is wrong}.
  \item \textbf{Code} answers are runnable with the usual imports (\texttt{pandas as pd}, \texttt{numpy as np}, \texttt{scipy.stats}, \texttt{sklearn}, \texttt{bs4}); often more than one correct form exists --- the marker rewards any correct idiom.
  \item Work each question yourself first, then check here. The mark for each item matches the paper.
  \item Total: \textbf{100 marks}, 5 parts of 20.
\end{itemize}
\end{frame}

% =====================================================================
\section{Part A --- Python foundations}
% =====================================================================

\begin{frame}[fragile]{A1 --- \texttt{int(-3.9)} \hfill [2, MCQ]}
\textbf{Correct: B) \texttt{-3}.}
\begin{itemize}
  \item \texttt{int()} on a float \textbf{truncates toward zero} --- it drops the fractional part, it does \emph{not} round. So \texttt{int(-3.9)} $\to$ \texttt{-3} and \texttt{int(3.9)} $\to$ \texttt{3}.
\end{itemize}
\textbf{Why the others are wrong:}
\begin{itemize}
  \item A) \texttt{-4} would be \texttt{math.floor(-3.9)} (round toward $-\infty$) or \texttt{round(-3.9)} --- not what \texttt{int()} does.
  \item C) \texttt{3} confuses sign; truncation keeps the sign.
  \item D) No error --- \texttt{int()} of a float is perfectly valid.
\end{itemize}
\end{frame}

\begin{frame}[fragile]{A2 --- Mutable default argument \hfill [2, MCQ]}
\textbf{Correct: B) \texttt{['mug', 'pen']}.}
\begin{itemize}
  \item The default \texttt{cart=[]} is created \textbf{once, at definition time}, and \emph{shared} across all calls that don't pass their own list. The first call appends \texttt{"mug"}; the same list survives, so the second call appends \texttt{"pen"} to it.
\end{itemize}
\textbf{Fix (the \texttt{None}-sentinel idiom):}
\begin{lstlisting}[language=Python]
def add_to_cart(item, cart=None):
    if cart is None:
        cart = []
    cart.append(item)
    return cart
\end{lstlisting}
\textbf{Distractors:} A/C assume a fresh list per call (the trap); D ignores the append.
\end{frame}

\begin{frame}[fragile]{A3 --- f-string formatting \hfill [3, short]}
\begin{lstlisting}[language=Python]
total = 1234.5
return_rate = 0.0873
print(f"\u20ac{total:,.2f}")      # -> the euro sign then 1,234.50
print(f"{return_rate:.1%}")       # -> 8.7%
\end{lstlisting}
\textbf{Key points:}
\begin{itemize}
  \item \texttt{:,.2f} = thousands separator \emph{and} 2 decimals $\Rightarrow$ \texttt{1,234.50}.
  \item \texttt{:.1\%} multiplies by 100 and appends \texttt{\%} with 1 decimal $\Rightarrow$ \texttt{8.7\%} (0.0873 $\to$ 8.73 $\to$ rounded 8.7).
  \item 1 mark for the currency/thousands format, 1 for the percent format, 1 for exact rounding/output.
\end{itemize}
\end{frame}

\begin{frame}[fragile]{A4 --- Counting with \texttt{Counter} \hfill [4, code]}
\begin{lstlisting}[language=Python]
from collections import Counter

categories = ["mug", "pen", "mug", "lamp", "mug", "pen"]
counts = Counter(categories)          # Counter({'mug': 3, 'pen': 2, 'lamp': 1})

top = max(counts, key=counts.get)     # 'mug'
print(counts)
print(top)
\end{lstlisting}
\textbf{Notes:}
\begin{itemize}
  \item \texttt{Counter(iterable)} tallies occurrences in one line.
  \item \texttt{max(counts, key=counts.get)} returns the key with the largest count; \texttt{counts.most\_common(1)[0][0]} is an equally valid alternative.
  \item A manual \texttt{d[k] = d.get(k, 0) + 1} loop also earns full marks.
\end{itemize}
\end{frame}

\begin{frame}[fragile]{A5 --- Debug: accumulator bug \hfill [4, debug]}
\textbf{Bug:} inside the loop the code \emph{overwrites} \texttt{total} with each price (\texttt{total = p}) instead of \emph{adding} to it, so it returns only the last price.
\begin{lstlisting}[language=Python]
def cart_total(prices):
    total = 0
    for p in prices:
        total += p        # was: total = p
    return total
\end{lstlisting}
\begin{itemize}
  \item The classic \texttt{=} vs \texttt{+=} accumulator mistake.
  \item Equivalent Pythonic one-liner: \texttt{return sum(prices)}.
  \item Marks: 2 for identifying the overwrite, 2 for the correct \texttt{+=} fix.
\end{itemize}
\end{frame}

\begin{frame}[fragile]{A6 --- \texttt{ShoppingCart} class \hfill [5, code]}
\begin{lstlisting}[language=Python]
class ShoppingCart:
    TAX_RATE = 0.19                      # class attribute, shared

    def __init__(self):
        self.items = []                  # fresh list per instance (no mutable default)

    def add_item(self, name, price, qty=1):
        self.items.append({"name": name, "price": price, "qty": qty})

    def subtotal(self):
        return sum(it["price"] * it["qty"] for it in self.items)

    def total(self):
        return self.subtotal() * (1 + self.TAX_RATE)
\end{lstlisting}
\textbf{Marks:} class attr [1]; \texttt{items} in \texttt{\_\_init\_\_} not as a default [1]; \texttt{add\_item} [1]; \texttt{subtotal} [1]; \texttt{total} reusing \texttt{subtotal} + \texttt{self.TAX\_RATE} [1].
\end{frame}

% =====================================================================
\section{Part B --- Data, pandas, viz \& statistics}
% =====================================================================

\begin{frame}[fragile]{B1 --- Boolean filtering \hfill [2, MCQ]}
\textbf{Correct: B) \texttt{orders[(orders.returned == 1) \& (orders.total > 100)]}.}
\begin{itemize}
  \item Combine boolean Series with \texttt{\&} (element-wise), and \textbf{wrap each condition in parentheses} because \texttt{\&} binds tighter than the comparison operators.
\end{itemize}
\textbf{Why the others are wrong:}
\begin{itemize}
  \item A) Python \texttt{and} on a Series raises \texttt{ValueError} (ambiguous truth value).
  \item C) Without parentheses, \texttt{1 \& orders.total} is evaluated first $\Rightarrow$ wrong result / error.
  \item D) A comma inside \texttt{[]} is not valid row-filter syntax.
\end{itemize}
\end{frame}

\begin{frame}{B2 --- \texttt{WHERE} vs \texttt{HAVING} \hfill [2, MCQ]}
\textbf{Correct: B) \texttt{HAVING}.}
\begin{itemize}
  \item \texttt{HAVING} filters \emph{groups} after \texttt{GROUP BY} + aggregation --- it can see aggregates like \texttt{SUM(total)}.
\end{itemize}
\textbf{Why the others are wrong:}
\begin{itemize}
  \item A) \texttt{WHERE} filters \emph{raw rows} before grouping and cannot reference aggregates.
  \item C) \texttt{GROUP BY} defines the groups; it does not filter them.
  \item D) \texttt{ORDER BY} only sorts the final result.
\end{itemize}
\emph{Logical execution order:} FROM $\to$ WHERE $\to$ GROUP BY $\to$ HAVING $\to$ SELECT $\to$ ORDER BY $\to$ LIMIT.
\end{frame}

\begin{frame}[fragile]{B3 --- Return rate per category \hfill [3, short]}
\begin{lstlisting}[language=Python]
orders.groupby("category")["returned"].mean().sort_values(ascending=False)
\end{lstlisting}
\textbf{Why this works:}
\begin{itemize}
  \item \texttt{returned} is a 0/1 column, so its \textbf{mean is exactly the fraction of returned orders} (the return rate) --- the workhorse trick for rates.
  \item \texttt{.groupby("category")} splits by category; \texttt{.sort\_values(ascending=False)} ranks highest first.
  \item Marks: groupby+column [1]; \texttt{.mean()} as the rate [1]; descending sort [1].
\end{itemize}
\end{frame}

\begin{frame}[fragile]{B4 --- Two-sample test \hfill [4, code]}
\begin{lstlisting}[language=Python]
from scipy import stats

result = stats.ttest_ind(express, standard, equal_var=False)  # Welch's t-test
print(f"t = {result.statistic:.3f}, p = {result.pvalue:.4f}")
\end{lstlisting}
\textbf{Notes:}
\begin{itemize}
  \item \textbf{Welch's two-sample t-test} (\texttt{equal\_var=False}) is correct because variances and sample sizes are unequal --- do \emph{not} assume equal variance.
  \item Read the p-value from \texttt{result.pvalue}.
  \item Marks: correct function [2]; \texttt{equal\_var=False} + naming Welch [1]; printing \texttt{pvalue} [1].
\end{itemize}
\end{frame}

\begin{frame}[fragile]{B5 --- SQL: revenue per category \hfill [4, code]}
\begin{lstlisting}[language=SQL]
SELECT category, SUM(total) AS revenue
FROM orders
GROUP BY category
HAVING SUM(total) > 10000
ORDER BY revenue DESC;
\end{lstlisting}
\textbf{Notes:}
\begin{itemize}
  \item \texttt{GROUP BY category} + \texttt{SUM(total)} aggregates revenue per category.
  \item The group filter must use \texttt{HAVING} (not \texttt{WHERE}).
  \item \texttt{ORDER BY revenue DESC} may reference the \texttt{SELECT} alias (\texttt{ORDER BY} runs after \texttt{SELECT}).
  \item Marks: GROUP BY + SUM [2]; HAVING filter [1]; ORDER BY DESC [1].
\end{itemize}
\end{frame}

\begin{frame}[fragile]{B6 --- Debug: aggregate in \texttt{WHERE} \hfill [3, debug]}
\textbf{Bug:} \texttt{WHERE} runs \emph{before} grouping/aggregation, so \texttt{SUM(total)} does not yet exist there --- the engine raises a ``misuse of aggregate'' error. Group filters belong in \texttt{HAVING}.
\begin{lstlisting}[language=SQL]
SELECT category, SUM(total) AS revenue
FROM orders
GROUP BY category
HAVING SUM(total) > 10000;   -- was: WHERE SUM(total) > 10000
\end{lstlisting}
\begin{itemize}
  \item Marks: identify that an aggregate cannot appear in \texttt{WHERE} [2]; correct \texttt{HAVING} rewrite [1].
\end{itemize}
\end{frame}

\begin{frame}{B7 --- Interpreting p = 0.03 \hfill [2, MCQ]}
\textbf{Correct: B).} At $\alpha = 0.05$ we have evidence the means differ, but the p-value alone says nothing about the \emph{size} or importance of the effect.
\begin{itemize}
  \item Always pair a p-value with an \textbf{effect size} (e.g.\ Cohen's $d$) and a \textbf{confidence interval}.
\end{itemize}
\textbf{Why the others are wrong:}
\begin{itemize}
  \item A) A p-value is not the probability the means are equal.
  \item C) Significance $\neq$ large/important --- with enough data a tiny difference becomes significant.
  \item D) ``97\% chance the result is real'' is a common but false restatement of $1-p$.
\end{itemize}
\end{frame}

% =====================================================================
\section{Part C --- ML, evaluation \& feature engineering}
% =====================================================================

\begin{frame}[fragile]{C1 --- \texttt{.ravel()} order \hfill [2, MCQ]}
\textbf{Correct: B) \texttt{tn, fp, fn, tp}.}
\begin{itemize}
  \item \texttt{confusion\_matrix} lays rows out as \emph{actual} (negative, positive) and columns as \emph{predicted} (negative, positive); flattening row-major gives \texttt{[tn, fp, fn, tp]}.
\end{itemize}
\textbf{Why the others are wrong:} A, C, D all scramble that fixed order --- unpacking into the wrong names silently corrupts every downstream precision/recall number.
\end{frame}

\begin{frame}{C2 --- Scaling before the split \hfill [2, MCQ]}
\textbf{Correct: B).} Fitting the scaler on the whole dataset lets \textbf{test-set statistics (mean/std) leak into training}, so the reported score is optimistically inflated and won't hold in production.
\begin{itemize}
  \item Fix: put the scaler \emph{inside} a \texttt{Pipeline} so it is fitted on the training rows of each split/fold only.
\end{itemize}
\textbf{Why the others are wrong:}
\begin{itemize}
  \item A) True that trees are scale-invariant, but that is not \emph{why} pre-split scaling is wrong here.
  \item C) Scaling doesn't change class balance.
  \item D) \texttt{StandardScaler} can be re-fitted; that's not the issue.
\end{itemize}
\end{frame}

\begin{frame}{C3 --- Precision \& recall by hand \hfill [3, short]}
Given \textbf{TP = 40, FP = 10, FN = 20, TN = 930}:
\[
\text{Precision} = \frac{TP}{TP+FP} = \frac{40}{40+10} = \frac{40}{50} = 0.80
\]
\[
\text{Recall} = \frac{TP}{TP+FN} = \frac{40}{40+20} = \frac{40}{60} \approx 0.667
\]
\begin{itemize}
  \item Precision = ``of the orders we \emph{flagged}, how many were truly returned?''
  \item Recall = ``of the orders that were truly returned, how many did we \emph{catch}?''
  \item Marks: precision formula + value [1.5]; recall formula + value [1.5]. TN is not used.
\end{itemize}
\end{frame}

\begin{frame}[fragile]{C4 --- Leakage-free pipeline \hfill [4, code]}
\begin{lstlisting}[language=Python]
from sklearn.model_selection import train_test_split
from sklearn.pipeline import Pipeline
from sklearn.preprocessing import StandardScaler
from sklearn.linear_model import LogisticRegression

X_tr, X_te, y_tr, y_te = train_test_split(
    X, y, test_size=0.20, random_state=42, stratify=y)

pipe = Pipeline([("scaler", StandardScaler()),
                 ("clf", LogisticRegression(max_iter=1000))])
pipe.fit(X_tr, y_tr)
y_pred = pipe.predict(X_te)
\end{lstlisting}
\textbf{Marks:} split with \texttt{test\_size=0.20} + \texttt{stratify=y} [1.5]; \texttt{Pipeline(StandardScaler, LogisticRegression)} [1.5]; fit on train + predict on test [1]. The Pipeline is what makes it leakage-free.
\end{frame}

\begin{frame}[fragile]{C5 --- Expected-value offer threshold \hfill [5, applied]}
\textbf{(a)} Break-even churn probability:
\[
p^{*} = \frac{C}{s \cdot \text{CLV}} = \frac{8}{0.25 \times 400} = \frac{8}{100} = 0.08
\]
\textbf{(b)} Expected value at $p = 0.10$:
\[
EV = p \cdot s \cdot \text{CLV} - C = 0.10 \times 0.25 \times 400 - 8 = 10 - 8 = \boldsymbol{+\text{ EUR }2}
\]
\textbf{(c)} \textbf{Yes, send the offer:} the predicted churn probability $0.10 > p^{*} = 0.08$, equivalently the expected value is positive ($+$ EUR 2).
\begin{itemize}
  \item Marks: $p^{*}$ [2]; $EV$ [2]; correct decision with justification [1]. Note the threshold is \emph{per customer} --- it scales with CLV.
\end{itemize}
\end{frame}

\begin{frame}[fragile]{C6 --- Debug: precision formula \hfill [4, debug]}
\textbf{Bug:} the code divides by \texttt{(tp + fn)} --- that is \textbf{recall}. Precision divides by \texttt{(tp + fp)}.
\begin{lstlisting}[language=Python]
tn, fp, fn, tp = confusion_matrix(y_test, y_pred).ravel()
precision = tp / (tp + fp)      # was: tp / (tp + fn)
\end{lstlisting}
\begin{itemize}
  \item Mnemonic: \emph{precision} is about what you \emph{predicted positive} (add the \emph{false positives}); \emph{recall} is about the \emph{actual positives} (add the \emph{false negatives}).
  \item Marks: recognise it computes recall / wrong denominator [2]; correct \texttt{tp/(tp+fp)} [2].
\end{itemize}
\end{frame}

% =====================================================================
\section{Part D --- AI engineering}
% =====================================================================

\begin{frame}{D1 --- Temperature for classification \hfill [2, MCQ]}
\textbf{Correct: C) \texttt{0.0}.}
\begin{itemize}
  \item \texttt{temperature=0.0} makes the model (near-)deterministic --- the right choice for classification, extraction and any task you want to be reproducible/testable.
\end{itemize}
\textbf{Why the others are wrong:}
\begin{itemize}
  \item A) \texttt{1.0} and B) \texttt{0.7} inject randomness, appropriate for creative rewriting/brainstorming, not labelling.
  \item D) False --- temperature very much affects output variability.
\end{itemize}
\end{frame}

\begin{frame}{D2 --- MCP primitives \hfill [2, MCQ]}
\textbf{Correct: A) tools, resources, prompts.}
\begin{itemize}
  \item MCP's three primitives: \textbf{tools} (model-controlled actions), \textbf{resources} (application-controlled data, read by URI), \textbf{prompts} (user-controlled templates, like slash commands).
\end{itemize}
\textbf{Why the others are wrong:}
\begin{itemize}
  \item B) host / client / server are the three \emph{roles}, not the primitives (a tempting mix-up).
  \item C) GET/POST/PUT are HTTP methods; MCP rides on JSON-RPC 2.0.
  \item D) faithfulness/relevancy/recall are Ragas RAG-evaluation metrics.
\end{itemize}
\end{frame}

\begin{frame}{D3 --- Cosine similarity \& L2 normalization \hfill [3, short]}
\textbf{(a)} L2-normalize so that \textbf{vector magnitude / document length cannot fake similarity} --- cosine should measure \emph{direction} (meaning), not length; otherwise a longer document scores artificially high.

\medskip
\textbf{(b)} For unit-length (L2-normalized) vectors, $|a| = |b| = 1$, so
\[
\cos(a,b) = \frac{a \cdot b}{|a|\,|b|} = a \cdot b
\]
cosine similarity reduces to the plain \textbf{dot product}. That is why top-$k$ retrieval over a normalized matrix collapses to a single matrix multiply (\texttt{docs @ q}).
\begin{itemize}
  \item Marks: length/magnitude reasoning [2]; ``dot product'' [1].
\end{itemize}
\end{frame}

\begin{frame}[fragile]{D4 --- Cosine + top-k retrieval \hfill [4, code]}
\begin{lstlisting}[language=Python]
import numpy as np

def cosine(a, b):
    return a @ b / (np.linalg.norm(a) * np.linalg.norm(b))

def top_k(q, docs, k):
    sims = np.array([cosine(q, d) for d in docs])
    return np.argsort(-sims)[:k]          # indices of the k most similar docs
\end{lstlisting}
\textbf{Notes:}
\begin{itemize}
  \item \texttt{a @ b} is the dot product; divide by the norms for cosine.
  \item \texttt{np.argsort(-sims)} sorts \emph{descending}; slice \texttt{[:k]} for the top $k$.
  \item Vectorised alternative if \texttt{docs} rows are L2-normalized: \texttt{np.argsort(-(docs @ q))[:k]}.
  \item Marks: \texttt{cosine} [2]; correct top-$k$ by descending similarity [2].
\end{itemize}
\end{frame}

\begin{frame}[fragile]{D5 --- Tool schema (JSON Schema) \hfill [3, code]}
\begin{lstlisting}[language=Python]
get_return_rate_tool = {
    "name": "get_return_rate",
    "description": "Return the return rate for a given product category.",
    "parameters": {
        "type": "object",
        "properties": {
            "category": {
                "type": "string",
                "description": "Product category, e.g. 'mugs'."
            }
        },
        "required": ["category"],
    },
}
\end{lstlisting}
\textbf{Marks:} \texttt{name} + \texttt{description} [1]; \texttt{parameters} as a JSON-Schema object with a \texttt{string} \texttt{category} property [1]; \texttt{"required": ["category"]} [1]. The \texttt{description} is what the model reads to decide when to call it.
\end{frame}

\begin{frame}[fragile]{D6 --- Debug: bare \texttt{json.loads} \hfill [3, debug]}
\textbf{Why it crashes:} real LLM output is not guaranteed to be valid JSON --- it may add markdown \texttt{```json} fences or trailing prose --- so \texttt{json.loads} raises \texttt{JSONDecodeError}. \textbf{Always wrap it in try/except.}
\begin{lstlisting}[language=Python]
def classify(review):
    r = llm.chat(messages=[{"role": "user", "content": review}])
    try:
        return json.loads(r["text"])
    except json.JSONDecodeError:
        return {"sentiment": "unknown", "_raw": r["text"]}
\end{lstlisting}
\begin{itemize}
  \item Marks: identify the unguarded parse / non-JSON output [1.5]; try/except with a structured fallback [1.5]. (Bonus-worthy: also set \texttt{temperature=0} and ask for ``JSON only''.)
\end{itemize}
\end{frame}

\begin{frame}{D7 --- Validation vs correctness \hfill [3, applied]}
\textbf{(a)} The validator only checks the \textbf{shape} of the value: \texttt{"ORD"} is a non-empty string of the right type, so it \emph{passes}. Validation cannot tell that it is a \emph{truncated}, wrong value --- it checks \emph{shape, not truth}.

\medskip
\textbf{(b)} To catch a value that is well-formed but wrong you need \textbf{ground truth} --- a hand-labelled \textbf{golden set} to measure field-level accuracy against. Only comparing to known-correct answers reveals that \texttt{"ORD"} $\neq$ \texttt{"ORD-5521"}.
\begin{itemize}
  \item Marks: ``shape not truth'' explanation [1.5]; ground truth / golden set [1.5].
\end{itemize}
\end{frame}

% =====================================================================
\section{Part E --- Applied \& production}
% =====================================================================

\begin{frame}{E1 --- Why not a random split for forecasting \hfill [2, MCQ]}
\textbf{Correct: B).} A shuffled split scatters days so \textbf{future dates land in training} --- the model peeks at the future, test error looks great, then collapses on genuinely forward-only production data.
\begin{itemize}
  \item Fix: split by time (\texttt{.iloc[:cut]} / \texttt{.iloc[cut:]}, or \texttt{TimeSeriesSplit}); every test point must come \emph{after} every training point.
\end{itemize}
\textbf{Why the others are wrong:} A) speed is irrelevant; C) it's not about data volume; D) leakage affects both additive and multiplicative series.
\end{frame}

\begin{frame}[fragile]{E2 --- Docker layer caching \hfill [2, MCQ]}
\textbf{Correct: B)} copy \texttt{requirements.txt} and run \texttt{pip install} \emph{before} copying the source.
\begin{lstlisting}[language=bash]
COPY requirements.txt .
RUN pip install -r requirements.txt
COPY src/ ./src/          # changes often -> only this layer rebuilds
\end{lstlisting}
\begin{itemize}
  \item A code-only change invalidates only the late \texttt{COPY src/} layer; \texttt{pip install} stays cached. Rule: least-likely-to-change instructions first.
\end{itemize}
\textbf{Why the others are wrong:} A) \texttt{COPY . .} first re-runs \texttt{pip install} on every code edit; C) \texttt{python:latest} is a non-reproducible moving tag; D) \texttt{CMD} order doesn't affect the install cache.
\end{frame}

\begin{frame}[fragile]{E3 --- MAE and MAPE \hfill [3, short]}
Actual \texttt{[100, 200, 50]}, forecast \texttt{[110, 180, 60]}; absolute errors \texttt{[10, 20, 10]}.
\[
\text{MAE} = \frac{10 + 20 + 10}{3} = \frac{40}{3} \approx \boldsymbol{13.33}
\]
\[
\text{MAPE} = \frac{1}{3}\!\left(\tfrac{10}{100} + \tfrac{20}{200} + \tfrac{10}{50}\right)\!\times 100
= \frac{0.1 + 0.1 + 0.2}{3}\times 100 \approx \boldsymbol{13.33\%}
\]
\begin{itemize}
  \item MAE is in demand units; MAPE is scale-free (\%). Marks: MAE [1.5]; MAPE [1.5]. (Recall MAPE explodes if any actual is 0.)
\end{itemize}
\end{frame}

\begin{frame}[fragile]{E4 --- TF-IDF + KMeans topics \hfill [4, code]}
\begin{lstlisting}[language=Python]
from sklearn.feature_extraction.text import TfidfVectorizer
from sklearn.cluster import KMeans

vec = TfidfVectorizer(stop_words="english")
X = vec.fit_transform(reviews)

km = KMeans(n_clusters=3, n_init=10, random_state=42)
labels = km.fit_predict(X)
print(labels)          # one cluster id per review
\end{lstlisting}
\textbf{Marks:} \texttt{TfidfVectorizer(stop\_words="english")} + \texttt{fit\_transform} [2]; \texttt{KMeans(n\_clusters=3, n\_init=10, random\_state=...)} [1]; \texttt{fit\_predict} producing per-review labels [1]. (Top terms per topic $=$ \texttt{np.argsort(km.cluster\_centers\_[c])[::-1][:n]}.)
\end{frame}

\begin{frame}[fragile]{E5 --- Scraping the catalogue \hfill [4, code]}
\begin{lstlisting}[language=Python]
from bs4 import BeautifulSoup

soup = BeautifulSoup(html, "html.parser")
products = []
for li in soup.find_all("li", class_="product"):
    name = li.find("h2", class_="name").get_text(strip=True)
    price = li.find("span", class_="price").get_text(strip=True)
    products.append((name, price))
\end{lstlisting}
\textbf{Marks:} \texttt{BeautifulSoup(html, "html.parser")} [1]; \texttt{find\_all("li", class\_="product")} [1]; extract name + price via \texttt{find(...).get\_text()} (or \texttt{.text}) [1]; build \texttt{(name, price)} tuples [1]. (See E6 for the robust version.)
\end{frame}

\begin{frame}[fragile]{E6 --- Debug: \texttt{None.text} \hfill [3, debug]}
\textbf{Why it crashes:} when \texttt{find(...)} matches nothing it returns \texttt{None}, and \texttt{None.text} raises \texttt{AttributeError}. Capture the node, test for \texttt{None}, \emph{then} read it.
\begin{lstlisting}[language=Python]
node = li.find("h2", class_="name")
name = node.get_text(strip=True) if node is not None else None
\end{lstlisting}
\begin{itemize}
  \item A resilient extractor returns an all-\texttt{None} row instead of crashing, so you can \emph{alert on empty rows} rather than get paged at 3 a.m.
  \item Marks: explain \texttt{find} returns \texttt{None} [1.5]; \texttt{None}-guarded rewrite [1.5].
\end{itemize}
\end{frame}

\begin{frame}{E7 --- Scaling before KMeans on RFM \hfill [2, applied]}
\textbf{Essential step: standardize the RFM features} (e.g.\ \texttt{StandardScaler}) before clustering.
\begin{itemize}
  \item KMeans uses Euclidean distance. \textbf{Monetary} has by far the largest magnitude, so without scaling it \emph{dominates the distance} and the clustering degenerates into ``sort customers by spend'' --- Recency and Frequency are effectively ignored.
  \item Scaling puts all three axes on comparable footing so segments reflect the full RFM shape (champions, at-risk big spenders, bargain newcomers, \dots).
  \item Marks: name scaling/standardization [1]; explain the largest-magnitude-feature-dominates failure [1].
\end{itemize}
\end{frame}

\begin{frame}{Closing}
\begin{block}{Mark map --- 100 total}
\begin{itemize}
  \item \textbf{Part A} 2+2+3+4+4+5 = 20 \quad \textbf{Part B} 2+2+3+4+4+3+2 = 20
  \item \textbf{Part C} 2+2+3+4+5+4 = 20 \quad \textbf{Part D} 2+2+3+4+3+3+3 = 20
  \item \textbf{Part E} 2+2+3+4+4+3+2 = 20
\end{itemize}
\end{block}
\begin{itemize}
  \item Recurring themes: rates as means of 0/1 columns; leakage $=$ information from the future / the answer; validation checks \emph{shape}, ground truth checks \emph{truth}; decisions need a \emph{costed threshold}, not just a probability.
  \item This is one of three parallel-form mock exams --- sit all three for full coverage.
\end{itemize}
\end{frame}

\end{document}
